\begin{abstract}
% Context (What is needed to understand the "need"?)
%Declarative mapping languages describe the mapping
%	\todo{"the mapping"? Is there only one?}
%of heterogeneous data
%	\todo{I don't think the point is that the data may be heterogeneous (which it may very well be), but that the types of sources of input data may be heterogeneous.}
%to knowledge graph (KG).
%	\todo{Wrong English. Either: "to a knowledge graph" or ``to knowledge graphs"}
%Different flavours of mapping languages exist and each of these mapping languages
%has diverse implementations.  
%% Need (Why something needed to be done at all?)
%However, due to the lack of language-agnostic
%	\todo{You focus on the language-agnostic aspect here. Yet, none of the things you mention in the rest of this sentence requires a language-agnostic approach. These things could all be done based on language-specific formalizations.}
%well-defined formalizations,
%there are discrepancies across the different implementations of the same
%mapping language, complications in verifying the correctness of the
%implementations, and difficulty in studying the expressiveness of the
%mapping languages.  
%% Task (What was undertaken to address the need? It’s here that you write ‘in this paper, we …’)
%% Object (What the present document does or covers)
%% Findings (What the work done yielded or revealed)
%In this paper, we provide a solid theoretical
%	\todo{I wouldn't say that our contribution is "theoretical".}
%foundation for current and future
%works on mapping heterogeneous data to KG. 
%The main contribution is the formal definitions and a set of algebraic 
%mapping operators.
%Additionally, to show the expressiveness of our formalism, we provide a 
%translation algorithm from RML to our formalism; as a consequence, RML 
%is formalized using our formalism. 
%Furthermore, we provide a few rewriting rules to show the possible
%optimizations with our formalism. 
% Conclusion (What the findings mean for the audience)
% Perspectives (What the future holds, beyond this work)
%The algebraic mapping operators offer a theoretical foundation for defining the
%operational semantics of both current and future mapping languages,
%promoting a more unified approach to constructing knowledge graphs from
%heterogeneous data sources.
%
% \keywords{Mapping Algebra \and Formalization \and RML \and Knowledge Graph Construction.}

Although they exist since more than ten years already, have attracted diverse implementations, and have been used successfully in a significant number of applications, declarative mapping languages for constructing
   %RDF-based
knowledge graphs from heterogeneous types of data sources still lack a solid formal foundation.
	%The informal nature of existing definitions
	This
%
	%has led to discrepancies between different implementations. Moreover, it makes it impossible to introduce implementation and optimization techniques that are provably correct, and
	makes it impossible to introduce implementation and optimization techniques that are provably correct and, in fact, has led to discrepancies between different implementations. Moreover,
it precludes
	%studying computational properties or properly comparing different languages or language fragments in terms of their expressive power.
	%properly comparing fundamental properties of different languages (e.g., expressive power).
	studying fundamental properties of different languages (e.g., expressive power).
To address this gap, this paper introduces a
	%(language-agnostic)
	language-agnostic
algebra
	%to capture definitions of mappings from heterogeneous types of data sources to knowledge graphs.
	for capturing mapping definitions.
As further contributions, we show that the popular mapping language RML can be translated into our algebra%
	~(%
	%,
%
	%which has the side effect of providing
	by which we also provide
	%which establishes
%
	%the first
	a
%
formal definition of the semantics of RML%
	%\ mappings%
%
	%!)
	)
	%,
and we prove several algebraic
	%equivalences which can be used as rewriting rules
	rewriting rules that can be used
to optimize mapping plans based on our algebra.

% Instead of focusing on a single specific mapping language, ..
\end{abstract}
